\documentclass[10pt]{article}

\usepackage[utf8]{inputenc}

\usepackage[T1]{fontenc}

\usepackage{amsmath}

\usepackage{amsfonts}

\usepackage{amssymb}

\usepackage[version=4]{mhchem}

\usepackage{stmaryrd}

\usepackage{mathrsfs}

\usepackage{bbold}

\usepackage{CJKutf8}

\usepackage{graphicx}

\usepackage[export]{adjustbox}

\graphicspath{ {./images/} }



\begin{document}

каждая точка $x$ данного пространства $X$ имела окрестность $O_{x}$, замыкание к-рой в пространстве $X$ есть бикомпактное подпространство. Всякое локально бикомшактное хаусдорфово пространство и только такое пространство может рассматриваться как открытое множество нек-рого бикомпакта $\bar{X}$, причем $\bar{X}$ получается из $X$ присоединением одной только точки $\bar{x}$, и топология бикомпакта $\bar{X}$ однозначно определена последним требованием и топологией пространства $X$; бикомпакт $\bar{X}$ наз. о д н о т о ч е ч н ы м, или а ле кс а вдровским, бикомпактным р а спи и р ви е м пространства $X$.



Важнейшим после условия бикомпактности условием типа компактности является у с ло в и е па р а ко мп а к т н о с т и, требующее, чтобы за всяким открытым покрытием $\alpha$ данного пространства $X$ следовало локально конечное открытое покрытие $\alpha^{\prime}$ (семейство множеств Т. п. наз. л о к а л ь н о к о н е ч ны м в нем, если каждая точка его обладает окрестностью, к-рая может иметь общие точки лишь с конечным числом множеств семейства). Здесь уже нельзя требовать, чтобы $\alpha^{\prime}$ содержалось в $\alpha$. Все метризуемые пространства являются паракомпактными хаусдорфовыми пространствами.



Лит.: [1] А л е к с а н д р о в П. С., Введение в теорию множеств и общую топологию, М., 1977; [2] Б у р б а к и Н., Обща я топология. Основные структуры, пер. с франц., М., М., $1966 ;[4]$ А л е к с а н и р о в П. С., П а с ы н к о в Б. А., Введение в теорию размерности..., М., 1973. П. С. Александров.



ТОПОЛОГИЧЕСКОЕ ТЕНЗӦРНОЕ П РОИЗВЕДЕНИЕ локально выпуклых пространств $E_{1}$ и $E_{2}-$ локально выпуклое пространство, обладающее свойством универсальности по отношению к заданным на $E_{1} \times E_{2}$ билинейным операторам с нек-рым условием непрерывности. Точнее, пусть $\mathscr{K}-$ нек-рый класс локально выпуклых пространств и для каждого $F \in \widetilde{\varkappa}$ задано подмножество $T(F)$ множества рездельно неирерывных билинейных операторов из $E_{1} \times E_{2}$ в $F$. Тогда Т. т. п. $E_{1}$ и $E_{2}$ (относительно класса $T(F)$ ) наз. локально выпуклое пространство $E_{1} \tilde{\otimes} E_{2} \in \mathscr{\mathcal { K }}$ вместе с оператором $B \in T\left(E_{1} \bar{\otimes} E_{2}\right)$, обладающее следующим свойством: для любого $S \in T(F), F \in \mathcal{K}$, существует единственный непрерывный линейный оператор $R: E_{1} \tilde{\otimes} E_{2} \rightarrow F \quad$ такой, что $R \circ B=S$. Таким образом, если ситуация позволяет говорить о функторе $T: \mathscr{K} \rightarrow$ $\rightarrow \operatorname{Set} s$, то $E_{1} \tilde{\otimes} E_{2}$ определено как представляющий объект этого функтора.



Во всех известных (1985) примерах $\mathscr{\varkappa}$ содержит поле комплексных чисел $\mathbb{C}$, a $T(\mathbb{C})$ содержит все билинейные функционалы вида $f \circ g: f \in E_{1}^{*}, \quad g \in E_{1}^{*}$, переводящие $(x, y)$ в $f(x) g(y)$. В этом случае, если Т. т. п. существует, то в $E_{1} \tilde{\otimes} E_{2}$ есть плотное подпространство, к-рое можно отождествить с пространством $E_{1} \otimes E_{2}$ алгебраического тензорного произвеения; при этом $B(x, y)=x \otimes y$.



Если é состоит из всех раздельно (соответственно, совместно) непрерывных билинейных операторов, то Т. т. П. наз. и н д у к т и в н М (соответственно, п р о е к т и в н ы м). Наиболее важно проективное Т. т. п. Пусть $\left\{p_{i}\right\}$ - определяющие семейства полунорм в $E_{i}, i=1,2$; через $\pi$ обозначается топология в $E_{1} \otimes E_{2}$, определенная семейством полунорм



\[

\begin{aligned}

& \left\{p_{1} \widehat{\otimes} p_{2}\right\}: p_{1} \widehat{\otimes} p_{2}(u)= \\

& =\inf \left\{\sum_{k=1}^{n} p_{1}\left(x_{k}\right) p_{2}\left(x_{k}\right): \sum_{k=1}^{n} x_{k} \otimes y_{k}=u\right\} .

\end{aligned}

\]



Тогда если $\mathscr{K}$ - класс всех, соответственно, всех полных локально выпуклых пространств, то проективное Т. т. п. $E_{1}$ и $E_{2}$ существует и его локально выпуклое пространство есть $E_{1} \otimes E_{2}$ с топологией $\pi$, соответст- венно, его пополнение. Если $E_{i}$ - банаховы пространства с нормами $p_{i}, i=1,2$, то $p_{1} \widehat{\widehat{\otimes}} p_{2}$ - норма в $E_{1} \otimes E_{2}$, пополнение по к-рой обозначается через $E_{1} \widehat{\otimes} E_{2}$. Элементы $E_{1} \otimes E_{2}$ имеют для каждого $\varepsilon>0$ представление



\[

u=\sum_{k=1}^{\infty} x_{k} \otimes y_{k}

\]



где



\[

\sum_{k=1}^{\infty} p_{1}\left(x_{k}\right) p_{2}\left(y_{k}\right) \leqslant p_{1} \hat{\otimes} p_{2}(u)+\varepsilon .

\]



Если снабдить $E_{1} \bigotimes E_{2}$ более слабой, чем $\pi$, топологией с помощью семейства полунорм $p_{1} \otimes p_{2}$ :



\[

p_{1} \otimes p_{2}(u)=\sup _{f, g \in V \times W}|(f \otimes g) \stackrel{(u)}{V}|

\]



где $V$ и $W$ - поляры единичных шаров относительно $p_{1}$ и $p_{2}$, то возникает Т. т. п., иногда наз. слабым. Локально выпуклые пространства $E_{1}$, обладающие тем свойством, что для любого $E_{2}$ обе топологии в $E_{1} \bigotimes E_{2}$ совпадают, образуют важный класс ядерных пространств.



Проективное Т. т. П. связано с понятием с в о йс т в а а п п р о к с и м а ц и и: локально выпуклое пространство $E_{1}$ обладает свойством аппроксимации, если для каждого предкомпактного множества $K \subset E_{1}$ и окрестности нуля $U$ существует непрерывный оператор конечного ранга $\varphi: E_{1} \rightarrow E_{1}$ такой, что для всех $x \in K$ $x-\varphi(x) \subset U$. Все ядерные пространства обладают свойством аппроксимации. Банахово пространство $E_{1}$ обладает свойством аппроксимации тогда и только тогда, когда для любого банахова пространства $E_{2}$ оператор $\tau:\left[E_{1} \widehat{\otimes} E_{2}\right] \rightarrow\left(E_{\mathrm{I}}^{*} \otimes E_{2}^{*}\right)^{*}, \quad$ однозначно определенный равенством $[\tau(x \otimes y)](f \otimes g)=f(x) g(y)$, имеет нулевое ядро. Построено [3] сепарабельное банахово пространство без свойства аппроксимации (и тем самым доказано существование банаховых пространств без базиса Шаудера, поскольку последние всегда имеют свойство аппроксимации, - т. е. отрицательно решена т. н. «проблема базиса» С. Банаха).



Jlum.: [1] G r o t h e n d i e c k A., "Mem. Amer. Math. Soc.», 1955, №. 16, pt 1, p. $1-193$, pt 2, p. $1-140$; [2] WI e \begin{CJK}{UTF8}{mj}中\end{CJK} e p X., 1971; [3] E n f l o P., "Acta Math.», 1973, t., 130, № 3-4, p. 309-,



\includegraphics[max width=\textwidth, center]{2023_07_09_5a831f3cb04830abe7a0g-1}

ТопологИЯ - раздел математики, имеющий своим назначением выяснение и исследование, в рамках математики, идеи непрерывности. Интуитивно идея непрерывности выражает коренные свойства пространства и времени и имеет, следовательно, фундаментальное значение для познания. Соответственно, Т., в к-рой понятие непрерывности получает математич. воплощение, естественно вплетается почти во все разделы математики. В соединении с алгеброй Т. составляет общую основу математики и содействует ее единству.



Предметом топологии является исследование свойств фигур и их взаимного расположения, сохраняющихся гомеоморфизмами, т. е. взаимно однозначными и непрерывными в обе стороны отображениями. Следовательно, Т. можно квалифицировать как разновидность геометрии. Важной чертой этой геометрии является необычайная широта класса геометрич. объектов, попадающих в сферу действия ее законов.



Вызвана эта широта тем, что центральное понятие Т.- понятие гомеоморфизма не требует для своего определения никаких классич. геометрич. понятий типа расстояния, прямолинейности, линейности, гладкости и т. Д. Понятие гомеоморфизма и лежащее в его основе понятие непрерывного отображения предполагают только, что точки и множества точек рассматриваемой фигуры могут находиться в нек-ром интуитивно ясном отношении близости, отличном, вообще говоря, от простого отношения принадлежности.





\end{document}